\chapter{jump}
\lhead[tempest 2000]{}
\label{sec:listing}
\lstset{style=68KStyle}

Jumping off the web was a new game feature introduced by Tempest 2000. It allows the
player to pop off the web and avoid colliding with enemies. This behaviour is managed
by the routine responsible for updating the player's claw, \icode{claw\_con1}. We encountered
this routine briefly in
\hyperref[sec:updating]{\textcolor{blue}{'mainloop'}}.
Here is the part concerned with detecting if the player has pressed the 'Fire 2' button, which
intiates a jump.

\begin{lstlisting}
chek4jump:
        cmp #-2,wave_tim           ; Is the enemy wave over?
        beq cce                    ; If so, don't jump.
        tst jenable                ; Is jumping allowed?
        beq cce                    ; If not, don't jump.
        move.l fire_2,d0           ; Get the 'fire 2' button.
        and.l pad_now,d0           ; Check if it's been pressed.
        beq cce                    ; If not, skip to cce, otherwise..
        move.l #-$20000,cjump      ; .. start a jump
        move #6,sfx                ; Select the 'Player Jump' sound effect.
        jsr fox                    ; Make a v. silly noise
cce:    move.l pad_now,d0          ; Get currently pressed buttons.
        bra clawcon                ; Do the rest of the claw update routine.
\end{lstlisting}

From looking at the above, maybe you can divine that processing the jump is going to be
managed by a single variable called \icode{cjump}. All we do here is initialize it
to the slightly eye-watering value of \icode{-\$20000}, which is -131072 in decimal. As we
shall see, there is a trick to this number that allows the jump to be implemented nice
and compactly.

\clearpage
Before we see how the magic was achieved, let's briefly consider what is actually required
to achieve a jump-like effect. We visualize this in the simplified sequence of diagrams 
below: the claw moves gradually off the top of the web along its Z axis. That's it.

\begin{minipage}[c]{0.48\linewidth}
\begin{figure}[H]
    \centering
    \begin{adjustbox}{width=6cm,center}
      \includegraphics[width=12cm]{src/jump/web_29.png}%
    \end{adjustbox}
\end{figure}
\end{minipage}
\begin{minipage}[c]{0.06\linewidth}
\begin{figure}[H]
    \centering
    \begin{tikzpicture}[baseline={(current bounding box.south)}]
      \draw[->,line width=3pt] (0,0) to (1,0);
    \end{tikzpicture}
\end{figure}
\end{minipage}
\begin{minipage}[c]{0.48\linewidth}
\begin{figure}[H]
    \centering
    \begin{adjustbox}{width=6cm,center}
      \includegraphics[width=12cm]{src/jump/web_15.png}%
    \end{adjustbox}
\end{figure}
\end{minipage}

\begin{minipage}[c]{0.48\linewidth}
\begin{figure}[H]
    \centering
    \begin{adjustbox}{width=6cm,center}
      \includegraphics[width=12cm]{src/jump/web_15.png}%
    \end{adjustbox}
\end{figure}
\end{minipage}
\begin{minipage}[c]{0.06\linewidth}
\begin{figure}[H]
    \centering
    \begin{tikzpicture}[baseline={(current bounding box.south)}]
      \draw[->,line width=3pt] (0,0) to (1,0);
    \end{tikzpicture}
\end{figure}
\end{minipage}
\begin{minipage}[c]{0.48\linewidth}
\begin{figure}[H]
    \centering
    \begin{adjustbox}{width=6cm,center}
      \includegraphics[width=12cm]{src/jump/web_1.png}%
    \end{adjustbox}
\end{figure}
\end{minipage}

So to achieve a jump, we just have to adjust the Z position of the claw every time we
are updating its position in \icode{claw\_con1}. There is one added complication of course,
the claw has to come back down to earth again.

\begin{minipage}[c]{0.48\linewidth}
\begin{figure}[H]
    \centering
    \begin{adjustbox}{width=6cm,center}
      \includegraphics[width=12cm]{src/jump/web_1.png}%
    \end{adjustbox}
\end{figure}
\end{minipage}
\begin{minipage}[c]{0.06\linewidth}
\begin{figure}[H]
    \centering
    \begin{tikzpicture}[baseline={(current bounding box.south)}]
      \draw[->,line width=3pt] (0,0) to (1,0);
    \end{tikzpicture}
\end{figure}
\end{minipage}
\begin{minipage}[c]{0.48\linewidth}
\begin{figure}[H]
    \centering
    \begin{adjustbox}{width=6cm,center}
      \includegraphics[width=12cm]{src/jump/web_15.png}%
    \end{adjustbox}
\end{figure}
\end{minipage}

\begin{minipage}[c]{0.48\linewidth}
\begin{figure}[H]
    \centering
    \begin{adjustbox}{width=6cm,center}
      \includegraphics[width=12cm]{src/jump/web_15.png}%
    \end{adjustbox}
\end{figure}
\end{minipage}
\begin{minipage}[c]{0.06\linewidth}
\begin{figure}[H]
    \centering
    \begin{tikzpicture}[baseline={(current bounding box.south)}]
      \draw[->,line width=3pt] (0,0) to (1,0);
    \end{tikzpicture}
\end{figure}
\end{minipage}
\begin{minipage}[c]{0.48\linewidth}
\begin{figure}[H]
    \centering
    \begin{adjustbox}{width=6cm,center}
      \includegraphics[width=12cm]{src/jump/web_29.png}%
    \end{adjustbox}
\end{figure}
\end{minipage}


This is where the slightly magic value of \icode{-\$20000} starts to come in handy. It has the
property of being negative. In assembler arithmetic, negative values have a very useful property
if you keep adding positive values to them. At first they keep getting smaller, but when they reach
zero any future additions will mean they start getting bigger again. In other words, they cycle
around like a clock. So in our case if we keep adding some value to \icode{-\$20000} it will eventually
reach \icode{\$00000} and then start climbing to \icode{\$20000} again. 

This means that all we have to do to implement our jump is to choose a value to add to \icode{cjump} and
set that as the claw's Z position every time we're updating the claw object in \icode{claw\_con1}. At first
the claw will get nearer to us, but once \icode{cjump} becomes positive again it will start receding from us 
once more. In this case, the value we choose to add each time is \icode{\$c11} (3089 in decimal). This results
in approximately 42 steps on the way up and 42 on the way down.

\begin{lstlisting}
        ; Process a claw jump if necessary.
cc2:    tst t2k                    ; Are we playing T2K?
        beq cce                    ; If not, skip to clawcon.
        ; Process a claw jump.
        move.l cjump,d0            ; Store the jump var in d0.
        beq chek4jump              ; Skip if not currently jumping..
        add.l d0,12(a6)            ; Otherwise set the claw's new Z position.
        move.l 12(a6),d1           ; Store the claw's Z position in d1.
        move #webz-80,d2           ; Store the web's Z position in d2.
        swap d2                    ; Swap first 2 bytes with last 2 bytes.
        clr d2                     ; Clear the last 2 bytes of d2.
        sub.l d2,d1                ; Subtract the web's Z from the claw's Z.
        add.l vp_zbase,d1          ; Add the base viewpoint to the claw's Z.
        move.l d1,vp_z             ; Store as the new player's Z viewpoint.
        move.l d1,vp_ztarg         ; Store as the new player's Z viewpoint.
        add.l #$c11,cjump          ; Add 3089 to cjump.
\end{lstlisting}

Of course we also have to test for when the jump is over. We've 'hit the ground' when
the claw finally reaches the same Z position as the web again, or event slightly past it.
When this happens we clamp the claw to the web and reset everything so that the claw and
the player's viewpoint are back to normal again. That was easy.

\begin{lstlisting}
        tst.l d0                   ; Test prev val of cjump.
        bmi cce                    ; If less than 0, jump is still active.
        cmp #webz-80,12(a6)        ; If web Z less than claw Z..
        blt cce                    ; .. then jump is still active.
        move #webz-80,12(a6)       ; Otherwise jump over, attach claw to web.
        clr 14(a6)                 ; Clear the lower 2 bytes of claw's Z. 
        clr.l cjump                ; Clear the jump variable.
        and.l #$fffc0000,vp_z      ; Restore the player's viewpoint.
        and.l #$fffc0000,vp_ztarg  ; Restore the player's viewpoint.
\end{lstlisting}
